% -----------------------------------------------------------------------------
% SECTION: REFERENCES (AD, RD, TERMS & ABBREVIATIONS)
% -----------------------------------------------------------------------------
\section{References}

This section lists the normative Applicable Documents (\abbr{AD}) and informative Reference Documents (\abbr{RD}) governing the \abbr{EWAS} Space Segment (\abbr{SSEG}, Configuration Item \texttt{CI-03}, \abbr{WBS} Element \texttt{123456.03}). It also defines key project-specific technical terms and abbreviations used throughout this specification.

\subsection{Applicable Documents (AD)}

The following Applicable Documents contain normative provisions for the \abbr{EWAS} Space Segment. Through explicit reference in this specification, these documents establish binding requirements for system design, verification, operation, and disposal.

{\renewcommand{\arraystretch}{1.25}%
\arrayrulecolor{TableBorderBlue}%
\rowcolors{2}{white}{TableRowLightBlue}%
\noindent
\begin{tabularx}{\textwidth}{|>{\raggedright\arraybackslash}p{0.11\textwidth}|>{\raggedright\arraybackslash}p{0.26\textwidth}|>{\raggedright\arraybackslash}X|>{\raggedright\arraybackslash}p{0.11\textwidth}|}
    \hline
    \rowcolor{TableHeaderBlue}
    % COLUMN 1: AD No. & COLUMN 2: Document Number & COLUMN 3: Title         & COLUMN 4: Issue
    \thead{AD No.}     & \thead{Document Number}   & \thead{Title}           & \thead{Issue}     \\ \hline
    \textbf{AD-01}     & {EWAS-HSB-SYS-SRD-01}     & {EWAS Final Report: Mission Definition, CONOPS, and Overall System Requirements Baseline (R-001 to R-049)} & {Rev.\ 1.0}  \\ \hline
    \textbf{AD-02}     & {ECSS-E-ST-10C}           & {Space engineering --- System engineering general requirements} & {Rev.\ 1}  \\ \hline
    \textbf{AD-03}     & {ECSS-E-ST-10-06C}        & {Space engineering --- Technical requirements specification} & {Issue 1}  \\ \hline
    \textbf{AD-04}     & {ECSS-E-ST-10-02C}        & {Space engineering --- Verification} & {Rev.\ 1}  \\ \hline
    \textbf{AD-05}     & {ECSS-E-ST-10-03C}        & {Space engineering --- Testing} & {Issue 1}  \\ \hline
    \textbf{AD-06}     & {ECSS-E-ST-10-04C}        & {Space engineering --- Space environment} & {Rev.\ 1}  \\ \hline
    \textbf{AD-07}     & {ECSS-E-ST-20C}           & {Space engineering --- Electrical and electronic} & {Rev.\ 2}  \\ \hline
    \textbf{AD-08}     & {ECSS-E-ST-31C}           & {Space engineering --- Thermal control general requirements} & {Issue 1}  \\ \hline
    \textbf{AD-09}     & {ECSS-E-ST-32C}           & {Space engineering --- Structural general requirements} & {Rev.\ 1}  \\ \hline
    \textbf{AD-10}     & {ECSS-E-ST-50C}           & {Space engineering --- Communications} & {Issue 1}  \\ \hline
    \textbf{AD-11}     & {ECSS-Q-ST-30C / 40C}     & {Space product assurance --- Dependability and Safety} & {Rev.\ 1}  \\ \hline
    \textbf{AD-12}     & {ECSS-U-AS-10C}           & {Space sustainability --- Adoption Notice of ISO 24113 (Space debris mitigation requirements)} & {Rev.\ 1}  \\ \hline
    \textbf{AD-13}     & {ECSS-M-ST-10C}           & {Space project management --- Project planning and implementation} & {Rev.\ 1}  \\ \hline
    \textbf{AD-14}     & {ECSS-E-ST-10-24C}        & {Space engineering --- Interface management} & {Issue 1}  \\ \hline
\end{tabularx}}

\clearpage
\subsection{Reference Documents (RD)}

The following Reference Documents provide contextual engineering benchmarks, regulatory frameworks, and architectural guidelines consulted during the formulation of this specification.

{\renewcommand{\arraystretch}{1.25}%
\arrayrulecolor{TableBorderBlue}%
\rowcolors{2}{white}{TableRowLightBlue}%
\noindent
\begin{tabularx}{\textwidth}{|>{\raggedright\arraybackslash}p{0.11\textwidth}|>{\raggedright\arraybackslash}p{0.26\textwidth}|>{\raggedright\arraybackslash}X|>{\raggedright\arraybackslash}p{0.11\textwidth}|}
    \hline
    \rowcolor{TableHeaderBlue}
    % COLUMN 1: RD No. & COLUMN 2: Document Number & COLUMN 3: Title         & COLUMN 4: Issue
    \thead{RD No.}     & \thead{Document Number}   & \thead{Title}           & \thead{Issue}     \\ \hline
    \textbf{RD-01}     & {ESA-ADMIN-IPOL-2023}     & {ESA Space Debris Mitigation Policy and Zero Debris Charter} & {2023}  \\ \hline
    \textbf{RD-02}     & {CCSDS 131.0-B / 141.0-B} & {CCSDS Space Data System Standards for TM/TC and Optical Communications} & {Blue Book}  \\ \hline
    \textbf{RD-03}     & {ITU-RR-2024}             & {ITU Radio Regulations (Space Research, Earth Exploration-Satellite, and Inter-Satellite Services)} & {2024}  \\ \hline
    \textbf{RD-04}     & {EDA-ODINS-EYE-II}        & {European Defence Agency ODIN'S EYE II / TWISTER Early Warning Architecture Benchmarks} & {Ref.}  \\ \hline
    \textbf{RD-05}     & {SDA-PWSA-TRK}            & {Space Development Agency (SDA) Proliferated Warfighter Space Architecture --- Tracking Layer Tranche 1/2 Specifications} & {Ref.}  \\ \hline
    \textbf{RD-06}     & {ECSS-E-HB-10-02A}        & {Space engineering --- Verification guidelines} & {Issue 1}  \\ \hline
\end{tabularx}}

\vspace{1.0em}
\subsection{Terms, Definitions and Abbreviations}

For the purposes of this specification, the terms and definitions given in \textbf{ECSS-S-ST-00-01C} and \textbf{ECSS-E-ST-10-06C} (\textbf{AD-03}) apply. In addition, the following key project-specific technical terms and abbreviations are defined:
\begin{itemize}
    \item \textbf{Proton Despiking:} Real-time onboard filtering that deletes artificial bright pixel spots (transient ionization spikes) on the infrared detector caused by energetic protons in the space radiation environment.
    \item \textbf{South Atlantic Anomaly (\abbr{SAA}):} Geographic region over South America and the South Atlantic Ocean ($15^\circ\text{S}\text{--}45^\circ\text{S}$, $90^\circ\text{W}\text{--}0^\circ\text{E}$) where Earth's tilted ($\sim 11^\circ$) and offset ($\sim 500\text{ km}$) geomagnetic dipole causes the inner Van Allen proton belt to dip down to altitudes of $200\text{--}1{,}000\text{ km}$.
    \item \textbf{Careful \abbr{COTS}:} Parts-selection and radiation-assurance philosophy combining lot-screened, radiation-tolerant commercial/industrial electronics with localized graded-Z Aluminium/Tantalum ($\mathrm{Al}/\mathrm{Ta}$) spot shielding and fault-tolerant avionics (\abbr{EDAC}, \abbr{TMR}, and fast-trip latch-up current limiters).
\end{itemize}

{\renewcommand{\arraystretch}{1.20}%
\arrayrulecolor{TableBorderBlue}%
\rowcolors{2}{white}{TableRowLightBlue}%
\noindent
\begin{tabularx}{\textwidth}{|>{\raggedright\arraybackslash}p{3.2cm}|>{\raggedright\arraybackslash}X|}
    \hline
    \rowcolor{TableHeaderBlue}
    % COLUMN 1: Abbreviation          & COLUMN 2: Meaning
    \thead{Abbreviation}              & \thead{Meaning}                                   \\ \hline
    \textbf{ADCS}                     & Attitude Determination and Control Subsystem      \\ \hline
    \textbf{AIT}                      & Assembly, Integration, and Testing                \\ \hline
    \textbf{AKE}                      & Attitude Knowledge Error                          \\ \hline
    \textbf{APE}                      & Absolute Pointing Error                           \\ \hline
    \textbf{BOL}                      & Beginning of Life                                 \\ \hline
    \textbf{C\&DH}                    & Command and Data Handling                         \\ \hline
    \textbf{CCSDS}                    & Consultative Committee for Space Data Systems     \\ \hline
    \textbf{CI}                       & Configuration Item                                \\ \hline
    \textbf{COMSEC}                   & Communications Security                           \\ \hline
    \textbf{CONOPS}                   & Concept of Operations                             \\ \hline
    \textbf{COTS}                     & Commercial Off-The-Shelf                          \\ \hline
    \textbf{CVCM}                     & Collected Volatile Condensable Material           \\ \hline
    \textbf{DD}                       & Displacement Damage                               \\ \hline
    \textbf{ECSS}                     & European Cooperation for Space Standardization    \\ \hline
\end{tabularx}}

\newpage

{\renewcommand{\arraystretch}{1.12}%
\arrayrulecolor{TableBorderBlue}%
\rowcolors{2}{white}{TableRowLightBlue}%
\noindent
\begin{tabularx}{\textwidth}{|>{\raggedright\arraybackslash}p{3.2cm}|>{\raggedright\arraybackslash}X|}
    \hline
    \rowcolor{TableHeaderBlue}
    \thead{Abbreviation}              & \thead{Meaning (Continued)}                       \\ \hline
    \textbf{EDAC}                     & Error Detection and Correction                    \\ \hline
    \textbf{EOL}                      & End of Life                                       \\ \hline
    \textbf{EPS}                      & Electrical Power Subsystem                        \\ \hline
    \textbf{EWAS}                     & Early Warning Alarm System                        \\ \hline
    \textbf{FDIR}                     & Fault Detection, Isolation, and Recovery          \\ \hline
    \textbf{FOR}                      & Field of Regard                                   \\ \hline
    \textbf{FOV}                      & Field of View                                     \\ \hline
    \textbf{FPA}                      & Focal Plane Array                                 \\ \hline
    \textbf{FSM}                      & Fast Steering Mirror                              \\ \hline
    \textbf{GEO}                      & Geostationary Earth Orbit (35,786~km)             \\ \hline
    \textbf{GSD}                      & Ground Sampling Distance                          \\ \hline
    \textbf{HIL}                      & Hardware-in-the-Loop                              \\ \hline
    \textbf{IIP}                      & Instantaneous Impact Point                        \\ \hline
    \textbf{IR}                       & Infrared                                          \\ \hline
    \textbf{ISL}                      & Inter-Satellite Link                              \\ \hline
    \textbf{LEO}                      & Low Earth Orbit (1,000~km)                        \\ \hline
    \textbf{LEOP}                     & Launch and Early Orbit Phase                      \\ \hline
    \textbf{LOS}                      & Line of Sight                                     \\ \hline
    \textbf{MCT}                      & Mercury Cadmium Telluride ($\mathrm{HgCdTe}$)     \\ \hline
    \textbf{MLI}                      & Multi-Layer Insulation                            \\ \hline
    \textbf{MMOD}                     & Micrometeoroids and Orbital Debris                \\ \hline
    \textbf{MWIR}                     & Medium-Wave Infrared ($3.0\text{--}5.0~\mu\mathrm{m}$) \\ \hline
    \textbf{NETD}                     & Noise Equivalent Temperature Difference           \\ \hline
    \textbf{NFOV}                     & Narrow Field of View                              \\ \hline
    \textbf{OBDH}                     & On-Board Data Handling Subsystem                  \\ \hline
    \textbf{OISL}                     & Optical Inter-Satellite Link                      \\ \hline
    \textbf{PLD}                      & Mission Payload                                   \\ \hline
    \textbf{PROP}                     & Propulsion Subsystem                              \\ \hline
    \textbf{RPE}                      & Relative Pointing Error                           \\ \hline
    \textbf{SAA}                      & South Atlantic Anomaly                            \\ \hline
    \textbf{SATG}                     & GEO Satellite Configuration Item (\abbr{WBS} \texttt{123456.03.02}) \\ \hline
    \textbf{SATL}                     & LEO Satellite Configuration Item (\abbr{WBS} \texttt{123456.03.01}) \\ \hline
    \textbf{SEE}                      & Single-Event Effects                              \\ \hline
    \textbf{SEU}                      & Single-Event Upset                                \\ \hline
    \textbf{SSEG}                     & Space Segment Configuration Item (\abbr{WBS} \texttt{123456.03}) \\ \hline
    \textbf{STR/MEC}                  & Structures and Mechanisms Subsystem               \\ \hline
    \textbf{SWIR}                     & Short-Wave Infrared ($1.0\text{--}3.0~\mu\mathrm{m}$) \\ \hline
    \textbf{TCS}                      & Thermal Control Subsystem                         \\ \hline
    \textbf{TID}                      & Total Ionizing Dose                               \\ \hline
    \textbf{TML}                      & Total Mass Loss                                   \\ \hline
    \textbf{TRS}                      & Technical Requirements Specification              \\ \hline
    \textbf{TT\&C}                    & Telemetry, Tracking, and Command                  \\ \hline
    \textbf{VCD}                      & Verification Control Document                     \\ \hline
    \textbf{WBS}                      & Work Breakdown Structure                          \\ \hline
    \textbf{WFOV}                     & Wide Field of View                                \\ \hline
\end{tabularx}}
